\scanpage{76}
On voit directement que $\unsure{f'|\wp'}(z)$ est elliptique en y changeant $z$ en $z+\unsure{\text{deux}|2\omega_i|2\omega_1}$ \deleted{\unsure{périodes}}. De $f'(z)-f'(z+2\omega_i)=0$ on déduit $f(z)-f(z+2\omega_i)=\mathrm{c}^{\mathrm{te}}$, \unsure{dite|cette} constante est en particulier égale à $f(-\omega_i)-f(\omega_i)$, or $f$ est paire, donc la constante est nulle. \unsure{Par suite}

Il existe une fonction elliptique et une seule (de périodes $2\omega_1$ et $2\omega_2$) qui n'ait comme pôles que les $z\in\unsure{\mathfrak P|\mathfrak B}$, \deleted{\unsure{et ait pour} pôle double et partie dominante $\frac1{z^2}$.} \inserted{et telle que $\unsure{\wp|p}(z)-\frac1{z^2}$ soit holomorphe \inserted{et nulle} à l'origine.} On la note $\wp(u)$, et elle est donnée par l'expression de Mittag-Leffler \unsure{(1)|(2)}. Cette fonction est paire, \deleted{et $\wp(z)-\frac1{z^2}$ est nulle \unsure{à l'origine}} :
\[
\wp z=\frac1{z^2}+c_1z^2+c_2z^4+\ldots
\]
(L'unicité résulte immédiatement de a) ; le reste de la proposition est immédiat)

c) \uline{Toute fonction elliptique \inserted{paire} \inserted{$f$} n'ayant de pôles qu'aux points de $\mathfrak P$, est un polynôme en $\wp$ \deleted{\unsure{,}} et réciproquement. -- Ce polynôme est unique}

La réciproque est évidente, ainsi que l'unicité, $\wp$ étant transcendant sur $C$. -- $f$ aura un pôle d'ordre pair $2k$ au point $0$, nous procédons par récurrence sur $k$ (évident si $k=0$). Or, si \deleted{$f(z)\ \frac{C_k}{\unsure{z^{2k}|\wp(z)}}$} si $f(u)=\frac{C_k}{u^{2k}}+\ldots$, $f(u)-C_k\wp(u)^k$ sera \deleted{\unsure{sera} fonction elliptique ayant un pôle \unsure{plus}} une fonction elliptique paire, n'ayant \deleted{d'ordre $2(k-1)$ au point} $0$, de pôles que dans $\mathfrak P$, et $0$ pour pôle d'ordre au plus $2(k-1)$, ce qui achève la démonstration.

\uline{\unsure{Noter|Notons} que le degré de $P$, \unsure{si|où} $f=P(\wp)$, est $k$}

d) \uline{Toute fonction elliptique paire $f$ est une fonction rationnelle en $\wp$, et réciproquement}. -- \uline{unicité}

Soit $P$ le parallélogramme sous-tendu par \unsure{$f$|$\mathfrak P$}, soient $a_1\ldots a_k$ \deleted{et \unsure{répétés un nombre de}} \inserted{l'ensemble des} pôles de $f$ \inserted{dans $P$}, non dans $\mathfrak P$, \deleted{et on suppose}

% archive sheet number: 75
